\section*{Bab \ref{ch:b2:solid-liquid-diagrams} --- Diagram Biner Padat--Cair}

\begin{solution}{exo:b2:solid-liquid-diagrams:1}
Pada \qty{1300}{\degreeCelsius}, \omterm{def:b2:solid-liquid-diagrams:liquidus}{likuidus} berada pada 0.541 dan \omterm{def:b2:solid-liquid-diagrams:liquidus}{solidus}
pada 0.609. 0.40: hanya cair. 0.58: cair (0.541) dan \omterm{def:b2:solid-liquid-diagrams:solid-solution}{larutan padat}
(0.609). 0.70: hanya \omterm{def:b2:solid-liquid-diagrams:solid-solution}{larutan padat}.
\end{solution}

\begin{solution}{exo:b2:solid-liquid-diagrams:2}
Fraksi padatan $(0.58 - 0.541)/(0.609 - 0.541) = 0.57$: \qty{1.15}{mol}
padatan, \qty{0.85}{mol} cairan. Nikel: $1.15 \times 0.609 = \qty{0.70}{mol}$
dalam padatan, $0.85 \times 0.541 = \qty{0.46}{mol}$ dalam cairan; total
\qty{1.16}{mol} $= 2.00 \times 0.58$.
\end{solution}

\begin{solution}{exo:b2:solid-liquid-diagrams:3}
Pada tekanan tetap $v = n + 1 - \varphi$ (tanpa reaksi). (a)
$1 + 1 - 2 = 0$. (b) $2 + 1 - 2 = 1$. (c) Cairan dan dua \omterm{def:b2:solid-liquid-diagrams:solid-solution}{larutan padat}:
$2 + 1 - 3 = 0$.
\end{solution}

\begin{solution}{exo:b2:solid-liquid-diagrams:4}
Penurunan mulus sampai \qty{232}{\degreeCelsius}, plato mendatar selama
timah membeku, lalu penurunan mulus lagi, yang makin lambat di dekat suhu
kamar. Pada plato, cairan dan padatan berdampingan pada tekanan tetap,
\omterm{def:b2:equilibrium-shifts:variance}{varians} 0: kalor yang dibuang dipasok oleh kristalisasi, dan suhu tidak
dapat berubah sampai cairan terakhir membeku.
\end{solution}

\begin{solution}{exo:b2:solid-liquid-diagrams:5}
Benzena: $\ln x_B = -(9870/8.314)(1/269.6 - 1/278.64) = -0.1429$, $x_B =
0.867$. Naftalena: $\ln x_N = -(19\,100/8.314)(1/269.6 - 1/353.2) = -2.017$,
$x_N = 0.133$. Jumlahnya 1.000: \qty{269.6}{K}
($\qty{-3.5}{\degreeCelsius}$) adalah suhu eutektik, dan cairan eutektik
mengandung 0.133 naftalena.
\end{solution}

\begin{solution}{exo:b2:solid-liquid-diagrams:6}
Cairan mendingin sampai \omterm{def:b2:solid-liquid-diagrams:liquidus}{likuidus} di sisi timbal, tempat kristal-kristal
(Pb) muncul; cairan makin kaya timah, sampai eutektik pada
\qty{182.2}{\degreeCelsius} dan 62.1~\%. Di sana sisa cairan membeku
menjadi \omterm{def:b2:solid-liquid-diagrams:eutectic}{campuran eutektik}. Fraksi massa eutektik:
$(30 - 18.91)/(62.13 - 18.91) = 0.26$. Pada pendinginan lebih lanjut,
kelarutan timah dalam (Pb) berkurang (garis putus-putus): sedikit (Sn)
mengendap di dalam kristal-kristal (Pb).
\end{solution}

\begin{solution}{exo:b2:solid-liquid-diagrams:7}
23.3~\% massa NaCl: $23.3/76.7 = \qty{0.304}{kg}$ garam per kilogram air.
Pada \qty{-25}{\degreeCelsius}, di bawah eutektik, tidak ada cairan yang
dapat ada berapa pun perbandingannya: garam dan es tetap berdampingan
sebagai padatan.
\end{solution}

\begin{solution}{exo:b2:solid-liquid-diagrams:8}
$x_{\mathrm B} = 0.5$ terletak di antara E$_1$ dan senyawa itu: $\mathrm{AB_2}$
mengkristal lebih dahulu, cairan bergerak menuju E$_1$, lalu membeku di
sana menjadi eutektik A dan $\mathrm{AB_2}$; pada akhirnya, A +
$\mathrm{AB_2}$. $x_{\mathrm B} = 0.9$ terletak di antara E$_2$ dan B: B
mengkristal lebih dahulu, cairan mencapai E$_2$; pada akhirnya,
$\mathrm{AB_2}$ + B.
\end{solution}

\begin{solution}{exo:b2:solid-liquid-diagrams:9}
Di tepi belakang zona, cairan membeku menjadi padatan yang hanya
mengandung 0.78 kali fraksi mol tembaganya: tembaga yang ditolak tetap
berada dalam cairan, yang terus bergerak dan makin kaya, sehingga tembaga
tersapu ke ujung batang. Dengan rasio yang mendekati 1, setiap lintasan
hanya menyingkirkan sebagian kecil pengotor: diperlukan banyak lintasan,
masing-masing dimulai dari batang yang ditinggalkan lintasan sebelumnya.
\end{solution}

\begin{solution}{exo:b2:solid-liquid-diagrams:10}
Penurunannya seperti dalam bab ini. Untuk cairan encer, $\ln x_A \approx -x_B$,
$1/T -
1/T^* \approx -\Delta T/T^{*2}$, sehingga
$x_B = \Delta_{\text{fus}}H_A\Delta T/(RT^{*2})$; dengan
$x_B \approx n_B/n_A = bM_A$, $\Delta T = (RT^{*2}M_A/\Delta_{\text{fus}}H_A)\,b$.
Benzena: $K_f = 8.314 \times 278.64^2 \times 0.07811/9870 = \qty{5.11}{K.kg/mol}$.
\end{solution}

\begin{solution}{exo:b2:solid-liquid-diagrams:11}
Per \qty{100}{g}: $19.0/24.31 = \qty{0.782}{mol}$ Mg, $81.0/207.2 =
\qty{0.391}{mol}$ Pb; rasio $2.00$: $\ce{Mg2Pb}$.
\end{solution}

\begin{solution}{exo:b2:solid-liquid-diagrams:12}
Plato sebanding dengan massa cairan eutektik: paduan yang tidak diketahui
mengandung $120/300 = 0.40$ bagiannya. Di sisi timbal,
$w = 18.91 + 0.40 \times (62.13 -
18.91) = 36.2~\%$ timah; di sisi timah,
$w = 96.59 - 0.40 \times (96.59 - 62.13) =
82.8~\%$. Patahan pertama \omterm{def:b2:solid-liquid-diagrams:cooling-curve}{kurva
pendinginan} (lebih tinggi untuk paduan 36~\%, yang \omterm{def:b2:solid-liquid-diagrams:liquidus}{likuidusnya} berada di
sisi timbal yang curam) atau sebuah mikrograf (kristal primer (Pb) atau
kristal primer (Sn)) membedakan keduanya.
\end{solution}

\begin{solution}{pb:b2:solid-liquid-diagrams:1}
\textbf{1.} Timbal \qty{327.5}{\degreeCelsius}, timah \qty{231.9}{\degreeCelsius}.
\textbf{2.} Per \qty{100}{g}: $62.13/118.71 = \qty{0.5234}{mol}$ timah dan
$37.87/207.2 = \qty{0.1828}{mol}$ timbal; $x_{\mathrm{Sn}} = 0.5234/0.7062 =
0.741$.
\textbf{3.} Titik-titik (0, 327.5), (100, 231.9), (62.13, 182.2),
(18.91, 182.2), (96.59, 182.2).
\textbf{4.} Seperti pada gambar: cair; L + (Pb); L + (Sn); (Pb); (Sn);
(Pb) + (Sn).
\textbf{5.} L (62.1~\% Sn) $\longrightarrow$ (Pb) (18.9~\%) + (Sn) (96.6~\%).
\textbf{6.} $v = 2 + 1 - 3 = 0$: suhu dan ketiga komposisi sudah tertentu.
\textbf{7.} \omterm{def:b2:solid-liquid-diagrams:solid-solution}{Larutan padat} yang kaya timbal (Pb): timbal melarutkan timah,
sehingga padatan yang berada dalam kesetimbangan dengan cairan sudah
mengandung timah, dalam proporsi yang dibaca pada \omterm{def:b2:solid-liquid-diagrams:liquidus}{solidus}.
\textbf{8.} Cairan itu mengikuti \omterm{def:b2:solid-liquid-diagrams:liquidus}{likuidus} menuju eutektik, makin kaya timah,
sampai 62.13~\%.
\textbf{9.} Cairan 62.13~\% dan (Pb) 18.91~\% timah.
\textbf{10.} (Pb): $(62.13 - 40)/(62.13 - 18.91) = 0.512$; cairan 0.488.
\textbf{11.} (Pb) 18.91~\% dan (Sn) 96.59~\%; (Sn): $(40 - 18.91)/(96.59 -
18.91) = 0.271$; (Pb): 0.729.
\textbf{12.} Konstituen eutektik 48.8~\% (cairan yang ada tepat di atasnya);
(Pb) primer 51.2~\%. (Pb) pada pertanyaan 11 adalah jumlah kristal primer
dan lamela (Pb) dalam eutektik.
\textbf{13.} Patahan pada \omterm{def:b2:solid-liquid-diagrams:liquidus}{likuidus}, penurunan yang lebih lambat, lalu plato
pada \qty{182.2}{\degreeCelsius} yang lamanya 0.488 kali plato paduan
eutektik.
\textbf{14.} $\ln x_{\mathrm{Sn}} = -(7030/8.314)(1/T - 1/505.1)$.
\textbf{15.} $\ln x_{\mathrm{Bi}} = -(11\,300/8.314)(1/T - 1/544.6)$.
\textbf{16.} $x_{\mathrm{Sn}}(T) + x_{\mathrm{Bi}}(T) = 1$.
\textbf{17.} Pada \qty{110}{\degreeCelsius}: $0.587 + 0.349 = 0.936$; pada
\qty{120}{\degreeCelsius}: $0.621 + 0.382 = 1.003$; pada
\qty{130}{\degreeCelsius}: $0.655 + 0.417 = 1.072$. Jumlahnya mencapai 1
di dekat \qty{119.5}{\degreeCelsius}: \qty{120}{\degreeCelsius} sampai
derajat terdekat.
\textbf{18.} $x_{\mathrm{Bi}} = 0.38$; dalam massa $0.38 \times 208.98/(0.38 \times
208.98 + 0.62 \times 118.71) = 0.52$.
\textbf{19.} Model memberikan \qty{120}{\degreeCelsius} dan 52~\% bismut,
diagram terkaji \qty{138.8}{\degreeCelsius} dan 57~\%: sekitar
\qty{19}{\degreeCelsius} terlalu rendah.
\textbf{20.} Timah yang mengkristal bukan timah murni, melainkan larutan
yang mengandung bismut: \omterm{def:b2:chemical-potential:chemical-potential}{potensial kimianya} lebih rendah daripada \omterm{def:b2:chemical-potential:chemical-potential}{potensial
kimia} timah murni, padatannya lebih stabil, dan padatan itu mengkristal
dari cairan pada suhu yang lebih tinggi untuk komposisi cairan tertentu.
Cabang timah \omterm{def:b2:solid-liquid-diagrams:liquidus}{likuidus} terangkat, dan bertemu dengan cabang bismut lebih
tinggi (cairan yang tidak ideal menambahkan koreksinya sendiri).
\textbf{21.} Solder itu melebur dan membeku pada satu suhu, yang terendah
dalam sistem itu: tidak ada rentang lembek yang dapat membuat sambungan
yang tergeser menjadi retak.
\textbf{22.} Timbal beracun, dan limbah elektronik berakhir di lingkungan.
\textbf{23.} Komponen yang peka panas dapat disolder pada suhu yang lebih
rendah; tetapi sambungannya melunak pada suhu yang lebih rendah selama
pemakaian, dan bismut membuatnya rapuh.
\textbf{24.} \qty{570}{g} bismut dan \qty{430}{g} timah.
\textbf{25.} Model cairan ideal meramalkan eutektik pada
$\boldsymbol{T_E
\approx \qty{120}{\degreeCelsius}}$, dibandingkan dengan
\qty{138.8}{\degreeCelsius} yang terkaji.
\end{solution}
